\begin{proposition}[The Izumi--Xu inversion obstruction]
\label{prop:m3-r436-unitary}
The based ring \(R8\_3\_436\) has no unitary categorification.
This statement does not exclude a general complex categorification.
\end{proposition}
\begin{proof}
The source indices \(\{0,1,2,5\}\) form the near-group ring
\(\operatorname{NG}(C_3,3)\), with \(g,g^2,X\) represented by source
indices \(1,2,5\), respectively, and
\[
 g^3=1,\qquad gX=Xg=X,\qquad X^2=1+g+g^2+3X.
\]
The invertible object with source index \(3\), denoted \(t\), satisfies
\(t^2=1\) and \(tgt^*=g^2\).  The independent crossed-ring construction
has model-to-source basis map \([0,3,1,7,2,6,5,4]\); the certificate checks
all \(512\) coefficients, as well as the displayed neutral subring and
conjugation action.

A unitary categorification would therefore contain a unitary
\(\operatorname{NG}(C_3,3)\) subcategory with a tensor autoequivalence
interchanging its nontrivial invertibles.  Izumi's C*-classification
\cite[Theorem~9.1 and Example~9.4]{IzumiNear}, together with
\cite[Example~13.7]{IzumiNear}, implies that every such unitary
near-group category has trivial tensor automorphism group.  This
contradicts conjugation by \(t\).

For the known Izumi--Xu categories, the same failure of the proposed
construction also follows from the triviality of their outer tensor
automorphism groups \cite[Section~4.2, following Remark~4.10, and
Remark~4.11]{GJSCyclic}: inner conjugation acts trivially on their
commutative based ring, whereas the required inversion does not.
The C*-classification does not classify all complex near-group
categories; the proposed nonunitary extension of its equations remains
conjectural in \cite[Remark~9.2]{IzumiNear}.  Thus the present argument
leaves the general complex categorification problem for this ring open.
\end{proof}
